% Input after finite_regularization.tex and before mean_field.tex.
% Conventional proof; labels use co:. No later-section dependencies.

\section{A one-parameter reduction for coherent ordering}
\label{co:section}

The four-row symmetry also permits an exact ordering theorem for a
single coherent evaluation. This theorem concerns the actual exterior
polynomial, including its correction terms.

Take \(m\) triples \((a_j,b_j,c_j)\) of nonnegative real numbers,
with \(a_j^2+b_j^2>0\) for each triple, and assume that at least one
triple has \(a_jb_jc_j>0\). An order of the triples defines a word
of \(4m\) rows by starting with the empty word and successively
prepending \((-a_j,b_j,c_j),(-a_j,-b_j,c_j)\) and appending
\((a_j,b_j,c_j),(a_j,-b_j,c_j)\). Indices \(j=0,\ldots,m-1\)
refer to this inner-to-outer order. For its rows \(v_i\), set
\[
 f(z)=\prod_i(v_i\cdot z),\qquad
 G_{ab}(z)=\sum_{i<j}(v_{ia}v_{jb}-v_{ib}v_{ja})
                         \prod_{h\ne i,j}(v_h\cdot z),
 \quad 0\le a<b\le2.
\]

Fix \(s>0\), put \(z_s=(1,is,0)\), and define
\begin{align*}
 S_j&=a_j^2+s^2b_j^2,&
 t_j&=\frac{a_jc_j}{S_j},&
 d_j&=\frac{a_j^2-s^2b_j^2}{S_j},&
 K_s&=\sum_j\frac{b_jc_j}{S_j},\\
 A&=\sum_j(2j+1)t_j,&
 B&=\sum_jt_j\left(d_j+2\sum_{\ell<j}d_\ell\right).
\end{align*}
The quantities \(A,B\) depend on the order; \(K_s\) does not.

\begin{proposition}[Exact coherent evaluation]\label{co:evaluation}
One has \(f(z_s)=\prod_jS_j^2>0\), and
\begin{equation}\label{co:components}
 \frac{G_{01}(z_s)}{f(z_s)}=0,\qquad
 \frac{G_{02}(z_s)}{f(z_s)}=4(A+B),\qquad
 \frac{G_{12}(z_s)}{f(z_s)}=4K_s+\frac{4i}{s}(B-A).
\end{equation}
Consequently
\begin{equation}\label{co:score}
 \mathcal E_s:=
 \frac{\sum_{a<b}|G_{ab}(z_s)|^2}{|f(z_s)|^2}
 =16\left((A+B)^2+\frac{(B-A)^2}{s^2}+K_s^2\right).
\end{equation}
At the unit point \(z_s/\sqrt{1+s^2}\), the corresponding ratio is
\((1+s^2)^2\mathcal E_s\).
\end{proposition}
\begin{proof}
The polynomial identities in the four-row recurrence hold for real
inputs as well as integers. Evaluate the squared variables at
\((X,Y,Z)=(1,-s^2,0)\). For the orbit factor \(p_j\),
\[
 p_j=S_j^2,\qquad
 \frac{(p_j)_X-s^2(p_j)_Y}{p_j}=2,\qquad
 \frac{(p_j)_X+s^2(p_j)_Y}{p_j}=2d_j.
\]
Thus, for the product \(F\) of the first \(j\) orbit factors,
\[
 \frac{F_X-s^2F_Y}{F}=2j,\qquad
 \frac{F_X+s^2F_Y}{F}=2\sum_{\ell<j}d_\ell.
\]
Substituting in \eqref{eq:rec02}--\eqref{eq:rec12} gives the
increments of the normalized ratios
\begin{align*}
 \Delta\frac{U_{02}-s^2U_{12}}F&=8(2j+1)t_j,\\
 \Delta\frac{U_{02}+s^2U_{12}}F
 &=8t_j\left(d_j+2\sum_{\ell<j}d_\ell\right).
\end{align*}
Equation~\eqref{eq:recV12} gives
\(\Delta(V_{12}/F)=4b_jc_j/S_j\).
Summing from the zero initial exterior polynomials and using
\(G_{02}=U_{02}\), \(G_{12}=xyU_{12}+V_{12}\), and
\(G_{01}=yzU_{01}+xzV_{01}\) proves \eqref{co:components}.
Its real and imaginary parts give \eqref{co:score}. Finally \(G\)
has degree two less than \(f\), which proves the unit-point scaling.
\end{proof}

\begin{theorem}[One-parameter sorting]\label{co:sorting}
For every fixed \(s>0\), an ordering maximizing \(\mathcal E_s\)
is obtained by sorting the triples in nondecreasing order of
\begin{equation}\label{co:keys}
 \frac{a_jc_j}{a_j^2+\kappa b_j^2}
\end{equation}
for some \(\kappa>0\), with suitable ordering among ties. The
family of sorts is independent of \(s\); its maximizing parameter
may depend on \(s\).
\end{theorem}
\begin{proof}
Since \(t_j\ge0\) and \(|d_j|\le1\), one has \(|B|\le A\).
The triple with all entries positive has \(t_j>0\) and
\(|d_j|<1\), so its contribution makes \(|B|<A\) in every order.
The order-independent term \(K_s^2\) can be dropped when maximizing.
Write
\[
 Q(A,B)=(A+B)^2+s^{-2}(B-A)^2
       =D(A^2+B^2)+2CAB,
 \qquad D=1+s^{-2},\quad C=1-s^{-2}.
\]
Here \(D>|C|\), so \(Q\) is positive definite.

Choose an order maximizing \(Q\), and denote its pair by \((A,B)\).
It also maximizes the linear functional \(A'+\lambda B'\), where
\begin{equation}\label{co:lambda}
 \lambda=\frac{CA+DB}{DA+CB}\in(-1,1).
\end{equation}
Indeed, \(DA+CB>0\), and a strict increase in this functional
would strictly increase \(Q\) by its quadratic expansion.
Moreover
\[
 (DA+CB)\mathbin\pm(CA+DB)=(D\pm C)(A\pm B)>0,
\]
which proves the strict parameter bounds in \eqref{co:lambda}.

If adjacent types \(i,j\) occur in that order, swapping them changes
\(A+\lambda B\) by
\[
 2\bigl((1+\lambda d_j)t_i-(1+\lambda d_i)t_j\bigr).
\]
Both factors \(1+\lambda d_j\) are positive. A strict inversion
in the keys \(t_j/(1+\lambda d_j)\) would therefore improve the
functional, contradicting maximality. The order is sorted in those
keys, and
\[
 \frac{t_j}{1+\lambda d_j}
 =\frac1{1+\lambda}
   \frac{a_jc_j}
        {a_j^2+s^2\frac{1-\lambda}{1+\lambda}b_j^2}.
\]
Taking \(\kappa=s^2(1-\lambda)/(1+\lambda)>0\) proves the theorem.
\end{proof}

\begin{corollary}[A finite sweep, including ties]\label{co:sweep}
At most \(\binom m2+1\) interval sorts suffice to attain the global
maximum in Theorem~\ref{co:sorting}. If the row coordinates and \(s^2\)
are rational, the sweep and the score comparisons can be carried out
with exact rational arithmetic.
\end{corollary}
\begin{proof}
Put \(u_j=a_jc_j\). Two keys that are not identical as functions
of \(\kappa\) cross at most once for \(\kappa>0\), at
\[
 \kappa=\frac{u_j a_i^2-u_i a_j^2}
                  {u_i b_j^2-u_j b_i^2},
\]
provided the denominator is nonzero and the quotient is positive.
These crossings partition \((0,\infty)\) into at most
\(\binom m2+1\) open intervals. Sort once at any parameter in
each interval.

To check that crossing ties cannot hide a larger value, let \(X=(A,B)\)
be a global maximizer and use its gradient functional from
\eqref{co:lambda}. Any other ordering \(Y\) tied for that functional
must have \(Y=X\): otherwise the quadratic expansion gives
\(Q(Y)=Q(X)+Q(Y-X)>Q(X)\).
There are finitely many orders. A sufficiently small one-sided
perturbation of the parameter cannot promote an order whose supporting
value was strictly smaller, so it selects a tied maximizer with the
same \((A,B)\). This covers simultaneous crossings as well.
Keys identical for all parameters have zero adjacent-swap increment
in both \(A\) and \(B\), so arbitrary tie-breaking within those
classes does not change the score. All crossing parameters, keys, and
values of \(Q\) are rational under the stated assumptions.
\end{proof}

Individual zero values of \(a_j\), \(b_j\), or \(c_j\) are permitted;
the assumption of one strictly positive triple ensures \(A>|B|\) and
keeps the supporting parameter interior. If all \(t_j=0\), then
\(A=B=0\) and \(\mathcal E_s=16K_s^2\) is order-independent; this
trivial case does not need a supporting parameter. Boundary-only
systems with \(A>0\) but no strictly positive triple are outside the
sorting theorem as stated.

The optimization proved here is at one specified complex point.
A larger value of \(\mathcal E_s\) does not by itself imply a larger
coefficient-square Fischer norm or a negative endpoint derivative.
The finite signs elsewhere in this paper are certified separately by
their exact integer Fischer-norm computations.
